\chapter{Research Methods} % Main chapter title

\label{Chapter5} % Change X to a consecutive number; for referencing this chapter elsewhere, use \ref{ChapterX}

\lhead{Chapter 5. \emph{Research Methods}} % Change X to a consecutive number; this is for the header on each page - perhaps a shortened title

A mixture of qualitative and quantitative research methods were used in this study: examination of literature, scoping interviews, a survey, follow-up interviews, and statistical analysis.
Figure \ref{fig_methodology} shows a flow chart of the methods that were used to meet each of the objectives listed in Chapter \ref{Chapter4}.
This chapter describes and justifies the employment of the selected research methods. % in the order of each objective.
% Figure \ref{contacts} lists some of the contacts gathered so far that could potentially be surveyed or interviewed.

\begin{figure}[htbp]
	\centering
	\includegraphics[width=\textwidth]{figures/Methodology03.eps}
	\rule{\textwidth}{0.5pt} % use line???
	\caption[Methodology]{Methods used to meet each of the objectives.}
	\label{fig_methodology}
\end{figure}


%----------------------------------------------------------------------------------------
%	SECTION 1
%----------------------------------------------------------------------------------------

\section{Literature Review}

Literature was reviewed in topics regarding collaboration and communication in design, BIM tools and interoperability.
Moreover, industry guidance documents in BIM adoption and work processes, specifically protocols, guides and scopes of services published by the industry and government, were closely examined.
% Hence, the principal research method employed here was the examination of literature, specifically protocols, guides and scopes of services published by the industry and government.

\section{Scoping Interviews}

Scoping interviews were conducted with two AEC professionals \citep{Quigley2017, Conaghan2017} in building services to identify:
\begin{itemize}
	\item Specific collaboration-related issues that have been experienced by building services engineers
	\item BIM tools used by building services engineers and their collaboration-enabling capabilities
	\item Authoritative industry documents that provide guidance for the adoption of BIM, e.g. PAS 1192-2 and ACE Schedule of Services MEP
\end{itemize}
\cite{Quigley2017} thought there was a need to clarify what consulting MEP engineers have to deliver information-wise at each stage of the process, and \cite{Conaghan2017} highlighted issues in consultant-to-contractor handovers.
Therefore, information regarding work stage deliverables and consultant-to-contractor handovers were sought in the literature.


\section{Survey} \label{sec_survey}

The purpose of the survey was to obtain subjective data from a group that might reveal the AEC industry's experience with collaborative work related to the building services.
A questionnaire was designed to verify whether the issues raised by \cite{Quigley2017} and \cite{Conaghan2017} resonated with others in industry, and to assess AEC professionals' familiarity and compliance with industry guidance.
Specifically, the issues that were to be verified regarded consultant-to-contractor handovers and work stage deliverables.

The questionnaire was designed using Google Forms and distributed via email or LinkedIn to the relevant AEC contacts of the author, their supervisor and Conaghan.
The contacts consisted mostly of consultants in building services, but also of contractors in building services, construction product manufacturers, policy and regulation consultants, and BIM managers or coordinators.
Policy and regulation consultants were contacted due to their potential to provide insight into the development of the industry guides.
The rest were contacted to provide insight into collaboration throughout the supply-chain (e.g. BIM management and consultant-to-contractor or consultant-to-specialist-designer handovers).
Additionally, an element of snowball sampling was used, where the author's, supervisor's and Conaghan's contacts were asked to forward the questionnaire to their colleagues.
The contacts had a window of approximately two weeks ($ 3^{rd} $ to $ 19^{th} $ February 2018) to complete and submit the questionnaire.

The questionnaire composed of an introduction to the research study and 29 questions divided among five sections: \say{Personal Information}, \say{Contractors Only}, \say{Work Stage Deliverables}, \say{Documents}, and \say{Thank you for completing this questionnaire!} (see Appendix \ref{Questionnaire}).
The majority of the questions provided quantitative data, through multiple-choice questions and two fully labelled five-point likert scales.
The scales were fully-labelled, as recommended by \cite{Vaerenbergh2012}, in order to avoid biased data and maximise credibility and validity when reporting direct summaries of responses, e.g. means and percentages.
The scales were used to ask the respondents to rate how much they agreed or disagreed with two statements.
Qualitative data was also collected in the form of five open-ended questions.
Two of these asked for additional comments or questions regarding Stage 4 BIM models and the questionnaire overall.
The other open-ended questions asked for explanations regarding the respondents' ratings on one of the scales, for examples from personal experience, and for email addresses if the respondents had agreed to participate in a follow-up interview or to a copy of the survey results.

\noindent
The purpose of the questions in the sections were to:
\begin{itemize}
	\item \say{Personal Information}:
	Bring context to the respondents’ answers by asking them about their profession, their work in building services, and their awareness of BIM.
	
	\item \say{Contractors Only}:
	Enquire contractors about BIM models at Stage 4 Technical Design of the RIBA PoW.
	
	\item \say{Work Stage Deliverables}:
	Enquire the respondents about their understanding of work stage deliverables and levels of definition, and to test the coherence of the industry's standard levels of detail.
	
	\item \say{Documents}:
	Assess the respondents' familiarity with some of the industry's main guidance documents.
	
	\item \say{Thank you for completing this questionnaire!}:
	Ask the respondents whether they would be willing to participate in a follow-up interview, whether they would like a copy of the final survey results, and to write any questions or comments, should they have any.
\end{itemize}
% \hl{Should this list be in Chapter 9 Industry Practice?}


\section{Follow-Up Interviews}

Statistical studies tend to benefit from large sample sizes, e.g. $ n_i>100 $.
% \hl{Cite Stats prof (Jenny?)...}
However, the sum of the number of people contacted was under 50.
Therefore, it was decided to conduct follow-up interviews with some of the respondents to provide more comprehensive and valid (though subjective) insights into the collaboration of practising building services engineers.


\section{Statistical Analysis}

The questionnaire data was entered into the IBM SPSS Statistics software programme to report direct summaries of the survey responses, e.g. means and percentages.
% Despite the small sample size…
% The five-point scale questions were set as ordinal level data.
% All the other questions were set as nominal level data.
% \hl{Double check this.}


\begin{comment}
\section{Collaboration and Communication in Design}

% Different kinds of literature were examined to provide context to the study and ascertain the guidance that building services engineers and the AEC industry overall receive regarding collaboration and work processes.
% In order to meet the first objective of exploring the meaning of collaboration in design, literature regarding the typical types of collaboration and communication in a design setting were sought.

\hl{Alex: content here about scoping interviews? semi-structured? in-person? how were interviewees selected?}

%----------------------------------------------------------------------------------------
%	SECTION 2
%----------------------------------------------------------------------------------------

\section{Tools and Interoperability}




%----------------------------------------------------------------------------------------
%	SECTION 3
%----------------------------------------------------------------------------------------

\section{Industry Guidance}

The third objective is to examine the industry's authoritative documents that provide guidance for building services engineers regarding collaboration and their work processes.
Hence, the principal research method employed here was the examination of literature, specifically protocols, guides and scopes of services published by the industry and government.
% Reasons for reading these documents, and by both government and industry?
Discussions with \cite{Quigley2017} and \cite{Conaghan2017} helped scope out what documents to examine (e.g. PAS 1192-2 and ACE Schedule of Services MEP) and what matters to look out for.
Specifically, \cite{Quigley2017} thought there was a need to clarify what consulting MEP engineers have to deliver information-wise at each stage of the process, and \cite{Conaghan2017} highlighted issues in consultant-to-contractor handovers.
Therefore, information regarding work stage deliverables and consultant-to-contractor handovers were sought in the literature.

Revelations from the survey (described in Section \ref{sec_ind_practice}) prompted further investigations of the industry guidance literature.
% expand???
% Overall, exclusively qualitative methods were  used to achieve the third objective.



%----------------------------------------------------------------------------------------
%	SECTION 4
%----------------------------------------------------------------------------------------

\section{Industry Practice} \label{sec_ind_practice}

The fourth objective is to assess the actual communication and collaboration processes of practising building services engineers.
Three research methods were employed for this: a survey, two follow-up interviews, and statistical analysis.

%------------------------------
%	SUBSECTION 1
%------------------------------

\subsection{Survey}

The purpose of the survey was to obtain subjective data from a group that might reveal the AEC industry's experience with collaborative work related to the building services.
A questionnaire was designed to verify whether the issues raised by \cite{Quigley2017} and \cite{Conaghan2017} resonated with others in industry, and to assess AEC professionals' familiarity and compliance with industry guidance.
Specifically, the issues that were to be verified regarded consultant-to-contractor handovers and work stage deliverables.

The questionnaire was designed using Google Forms and distributed via email or LinkedIn to the relevant AEC contacts of the author, their supervisor and Conaghan.
The contacts consisted mostly of consultants in building services, but also of contractors in building services, construction product manufacturers, policy and regulation consultants, and BIM managers or coordinators.
Policy and regulation consultants were contacted due to their potential to provide insight into the development of the industry guides.
The rest were contacted to provide insight into collaboration throughout the supply-chain (e.g. BIM management and consultant-to-contractor or consultant-to-specialist-designer handovers).
Additionally, an element of snowball sampling was used, where the author's, supervisor's and Conaghan's contacts were asked to forward the questionnaire to their colleagues.
The contacts had a window of approximately two weeks ($ 3^{rd} $ to $ 19^{th} $ February 2018) to complete and submit the questionnaire.

The questionnaire composed of an introduction to the research study and 29 questions divided among five sections: \say{Personal Information}, \say{Contractors Only}, \say{Work Stage Deliverables}, \say{Documents}, and \say{Thank you for completing this questionnaire!} (see Appendix \ref{Questionnaire}).
The majority of the questions provided quantitative data, through multiple-choice questions and two fully labelled five-point likert scales.
The scales were fully-labelled, as recommended by \cite{Vaerenbergh2012} in order to avoid biased data and maximise credibility and validity when reporting direct summaries of responses, e.g. means and percentages.
The scales were used to ask the respondents to rate how much they agreed or disagreed with two statements.
Qualitative data was also collected in the form of five open-ended questions.
Two of these asked for additional comments or questions regarding Stage 4 BIM models and the questionnaire overall.
The other open-ended questions asked for explanations regarding the respondents' ratings on one of the scales, for examples from personal experience, and for email addresses if the respondents had agreed to participate in a follow-up interview or to a copy of the survey results.

\noindent
The purpose of the questions in the sections were to:
\begin{itemize}
	\item \say{Personal Information}:
	Bring context to the respondents’ answers by asking them about their profession, their work in building services, and their awareness of BIM.
	
	\item \say{Contractors Only}:
	Enquire contractors about BIM models at Stage 4 Technical Design of the RIBA PoW.
	
	\item \say{Work Stage Deliverables}:
	Enquire the respondents about their understanding of work stage deliverables and levels of definition, and to test the coherence of the industry's standard levels of detail.
	
	\item \say{Documents}:
	Assess the respondents' familiarity with some of the industry's main guidance documents.
	
	\item \say{Thank you for completing this questionnaire!}:
	Ask the respondents whether they would be willing to participate in a follow-up interview, whether they would like a copy of the final survey results, and to write any questions or comments, should they have any.
\end{itemize}
% \hl{Should this list be in Chapter 9 Industry Practice?}


%------------------------------
%	SUBSECTION 2
%------------------------------

\subsection{Follow-Up Interviews}

Studies such as this tend to benefit from large sample sizes, e.g. $ n_i>100 $.
\hl{Cite Stats prof (Jenny?)...}
However, the sum of the number of people contacted was under 50 which reduced \hl{\ldots}.
Therefore, it was decided to conduct follow-up interviews with some of the respondents to provide more comprehensive and valid (though subjective) insight into collaboration in industry practice.
…

%------------------------------
%	SUBSECTION 3
%------------------------------

\subsection{Statistical Analysis}

Despite the small sample size…
The questionnaire data was entered into the IBM SPSS Statistics software programme.
The five-point scale questions were set as ordinal level data.
All the other questions were set as nominal level data.
\hl{Double check this.}
…

%----------------------------------------------------------------------------------------
%	SECTION 5
%----------------------------------------------------------------------------------------

\section{Comparative Study}


\end{comment}
